\begin{casebox}
\textbf{知识点小案例：同样读一百万个元素，地址顺序决定瓶颈。}
特例是三段循环：连续读 \code{a[i]}，按固定步长读 \code{a[i * 64]}，按随机索引读 \code{a[index[i]]}。三者逻辑上都是累加一百万个值，但机器看到的 cache line 利用率、TLB 覆盖和预取可能完全不同。连续读能把一条 cache line 的多个元素都用掉；大步长可能每次只用一个元素；随机索引还会让预取器失效。由此推出规律：复杂度相同不等于地址流相同。工程上分析内存问题时，先写出地址公式、工作集大小、页数和访问顺序，再看 cache/TLB 指标；不能只用 \code{O(n)} 解释内存性能。
\end{casebox}
